% 来源：OEBPS/Text/01_45.xhtml（OPF spine 顺序；仅第一册）
\cleardoublepage
\BookEntry[赫尔伯特·A.西蒙《我生活的种种模式》]{赫尔伯特·A.西蒙《我生活的种种模式》}

\BookCover{images/书单29我生活的种种模式.jpg}{[美] 赫尔伯特·A.西蒙著| 秦裕林等译| 东方出版中心2002年版}

\begin{BookSummaryBox}
\textbf{文理兼通的跨学科奇才，1978年诺贝尔经济学奖得主西蒙的自传}
\par
\end{BookSummaryBox}

\subsubsection*{这本书谈什么？}

这本书是1978年诺贝尔经济学奖得主赫尔伯特·A.西蒙的自传。西蒙是20世纪文理兼通的一位奇才，他从研究社会科学起步，在政治科学、经济学、心理学、管理科学、计算机科学和科学哲学等多个领域取得建树。从20世纪50年代开始，他把关于组织行为决策管理方面的理论应用于经济学，还开创了人工智能、认知科学和信息处理心理学等新学科，开辟了探索人类思维和智能的新天地。
\par

在这本自传中，西蒙用一个跨学科大学问家的眼光扫描了20世纪下半叶众多学科领域的发展，以及世界政治、经济局势的变化。这本书将告诉读者，现代科技与社会科学相结合后会对人类社会产生多么深远的影响。
\par

\subsubsection*{芒格为什么谈起这本书？}

1996年4月19日，时年72岁的查理·芒格在斯坦福大学法学院的演讲中，谈到了跨学科思维的重要性：
\par

\begin{quotation}
当我谈到跨学科方法——你们应该掌握各个学科的主要模型，并将它们统统派上用场——时，我是真的呼吁你们不要理会学科的法定界限。这个世界并不是按照跨学科的方法组织起来的。它反对跳出学科的法定范围。大企业也是这样……就这一点来说，我认为学术界错得有些离谱……如果你们的教授并没有教给你们正确的跨学科方法……你们可以自己改正那种愚蠢的做法。他是个笨蛋，并不意味着你们也要成为笨蛋。
\par
\end{quotation}

为了说明只有一种思维模型会让人变蠢，芒格最喜欢引用的一句谚语是：
\par

\begin{quotation}
在手里拿着铁锤的人看来，这个世界看上去就像一颗钉子。
\par
\end{quotation}

毫不奇怪，芒格会把跨学科奇才西蒙的这本自传推荐给他的读者。
\par

\subsubsection*{《我生活的种种模式》作者赫尔伯特·A.西蒙如是说}

“在大西洋北岸的美国，从威廉·詹姆斯那个时代开始，几乎一直到第二次世界大战为止，对人类思维的研究存在一个巨大的断层。美国心理学被行为科学主义、刺激与反应的联系、无意义音节和白鼠统治着。认知过程——耳朵在收到刺激到做出反应之间发生的事情——几乎没有被提及。‘心智’这个词留给了哲学家们，令人尊敬的心理学家是不讲这个词的。”
\par

“爱因斯坦在给哈达马特的一封信中写道：‘文字或语言（无论是书面的还是口头的）对我的思考机制似乎没有任何作用。那些似乎充当思维要素的心理实体是某种记号和或多或少清晰的形象。这种形象可以自动地再生和组合。’”
\par

“如果将科学看作是理论之间的竞争的话，它具有一个其他智力竞争所无法比拟的优点。在长期竞争中（不超过数个世纪），胜者的成功不是靠精湛的辩术，不是靠说服或迷惑外行听众的能力，也不是靠政治影响，而是依靠逐步积累的数据和事实的支持。只要事实的真实性不被推翻，人们就可以对一个理论的未来安之若素。”
\par
